% Source-aligned English for printed p.145, with logged indexing correction.
\SourcePage{145}

\LayerTarget{tohoku:section:2.4:intro}
\noindent spectral sequences we refer to [6, Chapters XV and XVII], confining
ourselves to specifying our terminology and isolating the most useful general
cases in which spectral sequences can be written down below.

\LayerTarget{tohoku:def:filtered-object}
Let $\mathbf{C}$ be an abelian category. A decreasing \textit{filtration} on an object
$A\in \mathbf{C}$ is a family $(F^n(A))$ of subobjects of $A$, indexed by
$n\in\mathbf Z$, such that $F^n(A)\subset F^{n'}(A)$ if $n\geqq n'$. A
\textit{filtered object} of $\mathbf{C}$ is an object equipped with a filtration. If
$A$ and $B$ are filtered objects, a morphism $u:A\to B$ is
compatible with the filtrations if $u(F^n(A))\subset F^n(B)$ for every $n$.
With these morphisms, the filtered objects of $\mathbf{C}$ form an additive category,
but not an abelian category, since a bijection in it need not be an
isomorphism.
\LayerTarget{tohoku:def:associated-graded}
The \textit{associated graded object} of a filtered object $A$ is the family
$G^n(A)=F^n(A)/F^{n+1}(A)$, denoted $G(A)$\textsuperscript{4)}\SourceFootnote{4}{These definitions are not
self-dual. Duality is restored by associating to every decreasing filtration
of $A$ by $F^n(A)$ the associated decreasing cofiltration
$F'_n(A)=A/F_{1-n}(A)$. On passage to the opposite category, these two become,
respectively, a decreasing cofiltration and its associated decreasing
filtration. It is also convenient to put $F_n(A)=F^{1-n}(A)$ and
$F'^n(A)=F'_{1-n}(A)$, the increasing filtration and cofiltration associated
to the preceding ones. Depending on the case, one of these four associated
filtrations is the most convenient. With
$G^n(A)=\mathrm{Coker.}(F^{n+1}(A)\to F^n(A))$ and
$G_n(A)=\mathrm{Ker}(F'_n(A)\to F'_{n-1}(A))$, one has
$G^n(A)=G_{n+1}(A)$. The functors $G^n$ and $G_n$ are interchanged on passage
to the opposite category.}. The
construction $G(A)$ is covariantly functorial
in the filtered object $A$.
\LayerTarget{tohoku:def:spectral-sequence}
A \textit{spectral sequence} in $\mathbf{C}$ is a system $E=(E_r^{p,q},E^n)$ consisting
of: a) objects $E_r^{p,q}\in \mathbf{C}$ for integers $p,q,r$ with $r\geqq2$; the 2
could be replaced by any integer $r_0$, although only $r=2$ or $r=1$ seems
interesting in applications; b) morphisms
$d_r^{p,q}:E_r^{p,q}\to E_r^{p+r,q-r+1}$
such that $d_r^{p+r,q-r+1}d_r^{p,q}=0$; c) isomorphisms
$\alpha_r^{p,q}:
\mathrm{Ker}(d_r^{p,q})/\mathrm{Im}(d_r^{p-r,q+r-1})
\to E_{r+1}^{p,q}$;
d) filtered objects $E^n$ of $\mathbf{C}$ for every integer $n$. For each fixed
$(p,q)$, one requires $d_r^{p,q}=0$ and $d_r^{p-r,q+r-1}=0$ for all
sufficiently large $r$. Thus $E_r^{p,q}$ is eventually independent of $r$;
the stable object is denoted $E_\infty^{p,q}$. For each fixed $n$, one also
requires $F^p(E^n)=E^n$ for sufficiently small $p$ and $F^p(E^n)=0$ for
sufficiently large $p$; and e) one is given isomorphisms
$\beta^{p,q}:E_\infty^{p,q}\to G^p(E^{p+q})$.
The family $(E^n)$, without its filtrations, is called the \textit{abutment}
of the spectral sequence $E$.
\LayerTarget{tohoku:def:morphism-spectral-sequence}
A \textit{morphism $u$ of spectral sequences} from $E=(E_r^{p,q},E^n)$ to
$E'=(E_r'^{p,q},E'^n)$ is a system of morphisms
$u_r^{p,q}:E_r^{p,q}\to E_r'^{p,q}$ and
$u^n:E^n\to E'^n$ compatible with the filtrations and commuting
with the $d_r^{p,q}$, $\alpha_r^{p,q}$, and $\beta^{p,q}$.
Spectral sequences in $\mathbf{C}$ form an additive category, though of course not an
abelian one.
\LayerTarget{tohoku:def:spectral-functor}
A \textit{spectral functor} is an additive functor from an abelian category to
a category of spectral sequences\textsuperscript{5)}\SourceFootnote{5}{For every known spectral functor,
the abutment seems in fact to be a cohomological functor. The relations between
the boundary homomorphisms in the abutment and the other constituents of the
spectral functor remain to be examined.}.
\LayerTarget{tohoku:def:cohomological-spectral-sequence}
A spectral sequence is \textit{cohomological} if $E_r^{p,q}=0$ for $p<0$ and
for $q<0$. Then $E_r^{p,q}=E_\infty^{p,q}$ once
$r>\mathrm{Sup}(p,q+1)$, $E^n=0$ for $n<0$, $F^m(E^n)=0$ if $m>n$, and
$F^m(E^n)=E^n$ if $m\leqq0$.

\TohokuContinuation{146}
